\documentclass[10pt,a4paper]{amsart}
\usepackage[margin=19mm]{geometry}
\usepackage{amsmath,amssymb,mathtools}
\usepackage[scr=rsfs,cal=euler]{mathalfa}
\usepackage{xfrac}
\usepackage[unicode,hidelinks,bookmarks=false]{hyperref}
\newcommand{\ie}{i.e.\ }
\newcommand{\cf}{cf.\ }
\newcommand{\resp}{respectively\ }
\newcommand{\Subsection}{\subsection}
\providecommand{\nobreakdash}{\nobreak\hbox{-}}
\setlength{\emergencystretch}{2em}
\multlinegap0pt
\usepackage{bm}
\usepackage{bbm}
\usepackage{dsfont}
\usepackage{xspace}
\usepackage{cancel}
\usepackage{mathtools}
\usepackage{amsthm}
\makeatletter
\@ifundefined{theorem}{\newtheorem{theorem}{Theorem}}{}
\@ifundefined{lemma}{\newtheorem{lemma}[theorem]{Lemma}}{}
\makeatother
\newcommand{\CGW}{\mathrm{CGW}}
\newcommand{\Cr}{\mathrm{C}}
\title[Complex generalised weighing matrices in centraliser algebra]{Complex generalised weighing matrices in centraliser algebras of monomial representations}
\author{Ronan Egan, Padraig \'O Cath\'ain, Andrea \v{S}vob}
\date{Source: arXiv:2607.16069v1 (CC BY 4.0)}
\begin{document}
\maketitle

\noindent\emph{Source and licence.} Complex generalised weighing matrices in centraliser algebras of monomial representations. Version arXiv:2607.16069v1 (CC BY 4.0), distributed under the Creative Commons Attribution 4.0 International licence, as recorded in the arXiv licence metadata for this exact version and shown on the abstract page https://arxiv.org/abs/2607.16069v1. Full rights notice: licenses/arxiv-2607.16069-CC-BY-4.0.txt.\par\smallskip

\noindent\textit{Editorial scope.} Counted unit: the Theorem whose source label is \texttt{thm:two-relation} in the article (source0.tex lines 527--533, source label \texttt{thm:two-relation}), delivered together with the article's own complete original proof of it (source0.tex lines 535--608, one \texttt{proof} environment of 74 lines). No mathematics is written, repaired or re-derived: statement and proof are the article's own text line for line. Required context retained uncounted: lem:basis-eigenvalues (lines 436--441); thm:hamming-flag (lines 496--502). Every definition, display and label of those lines is retained unchanged. Omitted from this package and not redistributed: 1--435, 442--495, 503--526, 534--534, 609--1001. Nothing inside a retained excerpt is omitted or altered: every retained body is delivered exactly as the article's lines are written, so no omission receipt of any kind accompanies this package. Citations: the retained excerpts carry no citation command at all, so no bibliography entry of the article is transcribed and no reference is redistributed. Reference adaptation: none was needed and none was made; every reference inside the delivered unit resolves inside it, so no unresolved reference or citation is produced. Printed slips kept literal: no printed word, symbol, hypothesis or number of the counted statement or of its proof was altered. Layout and class: the article's own class is \texttt{article} with packages including times, amssymb, amsmath, amsthm, float, multicol, caption, hyperref, fontenc, geometry, booktabs; the wrapper is the lane's pinned \texttt{amsart} editorial wrapper at 10pt on A4, which restates the theorem-like environments the retained text needs and adds the article's own macro definitions that the retained text needs (\texttt{\textbackslash CGW}, \texttt{\textbackslash Cr}), with extra wrapper packages (bm, bbm, dsfont, xspace, cancel, mathtools). The delivered numbering is the unit's own. Assets: no retained span contains a figure environment, a graphics-inclusion command, a PDF-page inclusion, a table or a TikZ picture; no image file is copied into this package, the package declares no asset dependency and no image file is redistributed with it. Licence: version arXiv:2607.16069v1 of this article is distributed under the Creative Commons Attribution 4.0 International licence, as recorded in the arXiv licence metadata for this exact version and shown on the abstract page https://arxiv.org/abs/2607.16069v1. A full rights notice is delivered at licenses/arxiv-2607.16069-CC-BY-4.0.txt. Characters: every retained excerpt is pure ASCII, so the delivered engine, XeLaTeX with its default Latin Modern text font, reports no missing character.
\begin{lemma}\label{lem:basis-eigenvalues}
If $Cv_l=\theta_l v_l$ for $l=1,\dots,d$, then $v_1\otimes\cdots\otimes v_d$ is an
eigenvector of $M_j$ with eigenvalue $e_j(\theta_1,\dots,\theta_d)$, the $j$-th
elementary symmetric polynomial in $\theta_1,\dots,\theta_d$.  As these products
range over the basis above, they give the entire spectrum of $M_j$.
\end{lemma}

\begin{theorem}\label{thm:hamming-flag}
Let $d\geq2$, let $\chi$ be a linear character of $\widehat H$, and let $M_j$
($1\leq j\leq d$) be a basis matrix of $\Cr(\rho_\chi)$.  Then $M_j$ is a weighing
matrix if and only if $j=d$ and the base $C$ is a conference matrix; this forces
$q$ even, and then $M_d=C^{\otimes d}\in\CGW(q^d,(q-1)^d;k)$.  In particular, for
$q$ odd the Hamming scheme $H(d,q)$ carries no flag-transitive weighing matrix.
\end{theorem}

\begin{theorem}\label{thm:two-relation}
Let $d=3$, and let $M$ be a sum of two basis matrices of $\Cr(\rho_\chi)$. Then $M \in \mathrm{CGW}(q^{3},w;k)$
only if the base $C$ is a conference matrix, hence $q$ even; the matrix $M$ is then
$I+S$, with $S=\beta\,C^{\otimes3}$ anti-Hermitian, for some $\beta \in \langle \zeta_{k} \rangle$, lying in
$\CGW\bigl(q^3,\,1+(q-1)^3;\,k\bigr)$, or it is one of two matrices on the binary
cube $H(3,2)$.
\end{theorem}

\begin{proof}
Scaling by a root of unity, take $M = M_a+\beta M_b$
with $0\le a<b\le3$.  Let $\sigma,\tau$ be the
eigenvalues of $C$ and $s=\tau/\sigma$; the relation $m\sigma+m'\tau=0$ makes
$s=-m/m'$ a negative rational.  By Lemma~\ref{lem:basis-eigenvalues} the common
eigenspaces are $V_0,\dots,V_3$, where $V_i$ is spanned by the tensors with
exactly $i$ factors an eigenvector for $\tau$; on $V_i$ the basis matrix $M_c$
($0\le c\le3$) acts as the scalar $\sigma^{c}Q_c(i)$, where $Q_c(i)$ is $e_c$
evaluated at $i$ copies of $s$ and $3-i$ copies of $1$ (the factor $\sigma^{c}$ is
pulled out by the homogeneity of $e_c$).  The values $Q_c(i)$ for $i=0,1,2,3$ are
\[
  Q_1:\ 3,\ s{+}2,\ 2s{+}1,\ 3s; \qquad
  Q_2:\ 3,\ 2s{+}1,\ s^2{+}2s,\ 3s^2; \qquad
  Q_3:\ 1,\ s,\ s^2,\ s^3.
\]
The eigenvalues of $M_{0}$ are all equal to $1$, so by abuse of notation we let $Q_{0}(i) = 1$ for all $i$. Put $p=\operatorname{Re}(\beta\sigma^{b-a})$ and $g=|\sigma|^{2(b-a)}$, so that
$p^2\le g$ (as $|\beta|=1$).  The eigenvalue of $M_a+\beta M_b$ on $V_i$ is
$\sigma^{a}Q_a(i)+\beta\sigma^{b}Q_b(i)$, of squared modulus
$|\sigma|^{2a}\bigl(Q_a(i)^2+2p\,Q_a(i)Q_b(i)+g\,Q_b(i)^2\bigr)$; hence the sum is
a weighing matrix precisely when
\begin{equation}\label{eq:wm-flat}
  Q_a(i)^2+2p\,Q_a(i)Q_b(i)+g\,Q_b(i)^2\ \text{ is independent of } i.
\end{equation}

First we show that if $M$ is a weighing matrix, we must have $s=-1$, so $C$ is a conference matrix and
$q$ is even.  If $a=0$ then 
\eqref{eq:wm-flat} requires the
quadratic $1+2pX+gX^2$ to take a single value at $X=Q_b(0),\dots,Q_b(3)$.  As
$g=|\sigma|^{2b}>0$ this quadratic is non-constant, so takes each value at most
twice, and the four numbers $Q_b(i)$ occupy at most two values.  For
$b=1$ they form an arithmetic progression of common difference $s-1\neq0$, hence
are distinct and no weighing matrix arises; for $b=2,3$ they pair up only at
$s=-1$.  If instead $1\le a<b$, regard the four instances $i=0,1,2,3$ of
\eqref{eq:wm-flat} as a linear system in $p$, $g$ and the common value $c$; a
solution exists only if
\[
  \det\bigl(Q_a(i)^2,\ Q_a(i)Q_b(i),\ Q_b(i)^2,\ 1\bigr)_{i=0}^{3}=0
\]
as each instance asserts that the row $\bigl(Q_a(i)^2,\,Q_a(i)Q_b(i),\,Q_b(i)^2,\,1\bigr)$
is orthogonal to the fixed nonzero vector $(1,\,2p,\,g,\,-c)$, so the four rows are
linearly dependent.  From the values above this determinant is a nonzero multiple of
$(s{-}1)^6(s{+}1)(s^2{+}4s{+}1)$, of $s(s{-}1)^6(s{+}1)^2(s^2{+}5s{+}1)$, and of
$s^3(s{-}1)^6(s{+}1)^3$ for $(a,b)=(1,2),(1,3),(2,3)$ respectively.  The factors
$s$ and $s-1$ and the quadratics $s^2+4s+1$, $s^2+5s+1$ (of discriminant $12$
and $21$) have no negative rational zero, so a negative rational $s$ annihilates
only the factor $s+1$; hence $s=-1$.  In every
non-vacuous case $s=-1$, i.e.\ $\tau=-\sigma$, so by the conference
characterisation above $C$ is a conference matrix and $q$ is even.

Finally we classify the solutions when $s=-1$.  Here $|\sigma|^2=q-1$, and the values
$Q_c(i)$ for $i=0,1,2,3$ become
\[
  Q_1:\ 3,1,-1,-3; \qquad Q_2:\ 3,-1,-1,3; \qquad Q_3:\ 1,-1,1,-1.
\]
For $\{M_0,M_3\}$ the values $Q_3(i)=\pm1$ make \eqref{eq:wm-flat} hold iff $p=0$,
with $g$ unconstrained: $\beta\sigma^3$ is imaginary, so $S:=\beta\,C^{\otimes3}$
is anti-Hermitian.  By Theorem~\ref{thm:hamming-flag},
$C^{\otimes3}\in\CGW(q^3,(q-1)^3;k)$, so $SS^*=(q-1)^3I$ and
\[
  (I+S)(I+S)^*=I+(S+S^*)+SS^*=\bigl(1+(q-1)^3\bigr)I
\]
for every even $q$ -- the stated family.  The other two-value pairs survive only
at $q=2$.  For $\{M_0,M_2\}$ the values $Q_2(i)\in\{3,-1\}$ give $p=-g$; with
$p^2\le g$ this forces $g\le1$, while $g=|\sigma|^4=(q-1)^2$, so $q=2$.  For
$\{M_1,M_3\}$ one finds $p=-1$ and $g=(q-1)^2$; here
$\beta\sigma^2/(q-1)=\beta\omega$ is a root of unity (using $\sigma^2=(q-1)\omega$)
with real part $p/(q-1)=-1/(q-1)$, and a root of unity with rational real part has
real part in $\{0,\pm\tfrac12,\pm1\}$, so $q-1\in\{1,2\}$ and, $q$ being even,
$q=2$.  The
remaining pairs are impossible: $\{M_1,M_2\}$ forces $g=-1$, and $\{M_2,M_3\}$
forces $p=0$ and $p=-2$ at once.  Hence for $q>2$ only the family $I+S$ occurs,
and at $q=2$ the further pairs $\{M_0,M_2\}$ and $\{M_1,M_3\}$ give two matrices
on the binary cube $H(3,2)$.
\end{proof}

\end{document}
